\documentclass[a4paper,12pt,twoside]{memoir}

% Castellano
\usepackage[spanish,es-tabla]{babel}
\selectlanguage{spanish}
\usepackage[utf8]{inputenc}
\usepackage[T1]{fontenc}
\usepackage{lmodern} % Scalable font
\usepackage{microtype}
\usepackage{placeins}

\RequirePackage{booktabs}
\RequirePackage[table]{xcolor}
\RequirePackage{xtab}
\RequirePackage{multirow}

% Links
\usepackage[colorlinks]{hyperref}
\hypersetup{
	allcolors = {red}
}

% Ecuaciones
\usepackage{amsmath}

% Rutas de fichero / paquete
\newcommand{\ruta}[1]{{\sffamily #1}}

% Párrafos
\nonzeroparskip


% Imagenes
\usepackage{graphicx}
\newcommand{\imagen}[2]{
	\begin{figure}[!h]
		\centering
		\includegraphics[width=0.9\textwidth]{#1}
		\caption{#2}\label{fig:#1}
	\end{figure}
	\FloatBarrier
}

\newcommand{\imagenflotante}[2]{
	\begin{figure}%[!h]
		\centering
		\includegraphics[width=0.9\textwidth]{#1}
		\caption{#2}\label{fig:#1}
	\end{figure}
}



% El comando \figura nos permite insertar figuras comodamente, y utilizando
% siempre el mismo formato. Los parametros son:
% 1 -> Porcentaje del ancho de página que ocupará la figura (de 0 a 1)
% 2 --> Fichero de la imagen
% 3 --> Texto a pie de imagen
% 4 --> Etiqueta (label) para referencias
% 5 --> Opciones que queramos pasarle al \includegraphics
% 6 --> Opciones de posicionamiento a pasarle a \begin{figure}
\newcommand{\figuraConPosicion}[6]{%
  \setlength{\anchoFloat}{#1\textwidth}%
  \addtolength{\anchoFloat}{-4\fboxsep}%
  \setlength{\anchoFigura}{\anchoFloat}%
  \begin{figure}[#6]
    \begin{center}%
      \Ovalbox{%
        \begin{minipage}{\anchoFloat}%
          \begin{center}%
            \includegraphics[width=\anchoFigura,#5]{#2}%
            \caption{#3}%
            \label{#4}%
          \end{center}%
        \end{minipage}
      }%
    \end{center}%
  \end{figure}%
}

%
% Comando para incluir imágenes en formato apaisado (sin marco).
\newcommand{\figuraApaisadaSinMarco}[5]{%
  \begin{figure}%
    \begin{center}%
    \includegraphics[angle=90,height=#1\textheight,#5]{#2}%
    \caption{#3}%
    \label{#4}%
    \end{center}%
  \end{figure}%
}
% Para las tablas
\newcommand{\otoprule}{\midrule [\heavyrulewidth]}
%
% Nuevo comando para tablas pequeñas (menos de una página).
\newcommand{\tablaSmall}[5]{%
 \begin{table}
  \begin{center}
   \rowcolors {2}{gray!35}{}
   \begin{tabular}{#2}
    \toprule
    #4
    \otoprule
    #5
    \bottomrule
   \end{tabular}
   \caption{#1}
   \label{tabla:#3}
  \end{center}
 \end{table}
}

%
% Nuevo comando para tablas pequeñas (menos de una página).
\newcommand{\tablaSmallSinColores}[5]{%
 \begin{table}[H]
  \begin{center}
   \begin{tabular}{#2}
    \toprule
    #4
    \otoprule
    #5
    \bottomrule
   \end{tabular}
   \caption{#1}
   \label{tabla:#3}
  \end{center}
 \end{table}
}

\newcommand{\tablaApaisadaSmall}[5]{%
\begin{landscape}
  \begin{table}
   \begin{center}
    \rowcolors {2}{gray!35}{}
    \begin{tabular}{#2}
     \toprule
     #4
     \otoprule
     #5
     \bottomrule
    \end{tabular}
    \caption{#1}
    \label{tabla:#3}
   \end{center}
  \end{table}
\end{landscape}
}

%
% Nuevo comando para tablas grandes con cabecera y filas alternas coloreadas en gris.
\newcommand{\tabla}[6]{%
  \begin{center}
    \tablefirsthead{
      \toprule
      #5
      \otoprule
    }
    \tablehead{
      \multicolumn{#3}{l}{\small\sl continúa desde la página anterior}\\
      \toprule
      #5
      \otoprule
    }
    \tabletail{
      \hline
      \multicolumn{#3}{r}{\small\sl continúa en la página siguiente}\\
    }
    \tablelasttail{
      \hline
    }
    \bottomcaption{#1}
    \rowcolors {2}{gray!35}{}
    \begin{xtabular}{#2}
      #6
      \bottomrule
    \end{xtabular}
    \label{tabla:#4}
  \end{center}
}

%
% Nuevo comando para tablas grandes con cabecera.
\newcommand{\tablaSinColores}[6]{%
  \begin{center}
    \tablefirsthead{
      \toprule
      #5
      \otoprule
    }
    \tablehead{
      \multicolumn{#3}{l}{\small\sl continúa desde la página anterior}\\
      \toprule
      #5
      \otoprule
    }
    \tabletail{
      \hline
      \multicolumn{#3}{r}{\small\sl continúa en la página siguiente}\\
    }
    \tablelasttail{
      \hline
    }
    \bottomcaption{#1}
    \begin{xtabular}{#2}
      #6
      \bottomrule
    \end{xtabular}
    \label{tabla:#4}
  \end{center}
}

%
% Nuevo comando para tablas grandes sin cabecera.
\newcommand{\tablaSinCabecera}[5]{%
  \begin{center}
    \tablefirsthead{
      \toprule
    }
    \tablehead{
      \multicolumn{#3}{l}{\small\sl continúa desde la página anterior}\\
      \hline
    }
    \tabletail{
      \hline
      \multicolumn{#3}{r}{\small\sl continúa en la página siguiente}\\
    }
    \tablelasttail{
      \hline
    }
    \bottomcaption{#1}
  \begin{xtabular}{#2}
    #5
   \bottomrule
  \end{xtabular}
  \label{tabla:#4}
  \end{center}
}



\definecolor{cgoLight}{HTML}{EEEEEE}
\definecolor{cgoExtralight}{HTML}{FFFFFF}

%
% Nuevo comando para tablas grandes sin cabecera.
\newcommand{\tablaSinCabeceraConBandas}[5]{%
  \begin{center}
    \tablefirsthead{
      \toprule
    }
    \tablehead{
      \multicolumn{#3}{l}{\small\sl continúa desde la página anterior}\\
      \hline
    }
    \tabletail{
      \hline
      \multicolumn{#3}{r}{\small\sl continúa en la página siguiente}\\
    }
    \tablelasttail{
      \hline
    }
    \bottomcaption{#1}
    \rowcolors[]{1}{cgoExtralight}{cgoLight}

  \begin{xtabular}{#2}
    #5
   \bottomrule
  \end{xtabular}
  \label{tabla:#4}
  \end{center}
}


















\graphicspath{ {./img/} }

% Capítulos
\chapterstyle{bianchi}
\newcommand{\capitulo}[2]{
	\setcounter{chapter}{#1}
	\setcounter{section}{0}
	\chapter*{#2}
	\addcontentsline{toc}{chapter}{#2}
	\markboth{#2}{#2}
}

% Apéndices
\renewcommand{\appendixname}{Apéndice}
\renewcommand*\cftappendixname{\appendixname}

\newcommand{\apendice}[1]{
	%\renewcommand{\thechapter}{A}
	\chapter{#1}
}

\renewcommand*\cftappendixname{\appendixname\ }

% Formato de portada
\makeatletter
\usepackage{xcolor}
\newcommand{\tutor}[1]{\def\@tutor{#1}}
\newcommand{\course}[1]{\def\@course{#1}}
\definecolor{cpardoBox}{HTML}{E6E6FF}
\def\maketitle{
  \null
  \thispagestyle{empty}
  % Cabecera ----------------
\noindent\includegraphics[width=\textwidth]{cabecera}\vspace{1cm}%
  \vfill
  % Título proyecto y escudo informática ----------------
  \colorbox{cpardoBox}{%
    \begin{minipage}{.8\textwidth}
      \vspace{.5cm}\Large
      \begin{center}
      \textbf{TFG del Grado en Ingeniería Informática}\vspace{.6cm}\\
      \textbf{\LARGE\@title{}}
      \end{center}
      \vspace{.2cm}
    \end{minipage}

  }%
  \hfill\begin{minipage}{.20\textwidth}
    \includegraphics[width=\textwidth]{escudoInfor}
  \end{minipage}
  \vfill
  % Datos de alumno, curso y tutores ------------------
  \begin{center}%
  {%
    \noindent\LARGE
    Presentado por \@author{}\\ 
    en Universidad de Burgos --- \@date{}\\
    Tutor: \@tutor{}\\
  }%
  \end{center}%
  \null
  \cleardoublepage
  }
\makeatother

\newcommand{\nombre}{Antonio de los Mozos Alonso} %%% cambio de comando

% Datos de portada
\title{Seguimiento de líneas basado en openCV para AGVs}
\author{\nombre}
\tutor{Jesús Enrique Sierra García}
\date{\today}

\begin{document}

\maketitle


\newpage\null\thispagestyle{empty}\newpage


%%%%%%%%%%%%%%%%%%%%%%%%%%%%%%%%%%%%%%%%%%%%%%%%%%%%%%%%%%%%%%%%%%%%%%%%%%%%%%%%%%%%%%%%
\thispagestyle{empty}


\noindent\includegraphics[width=\textwidth]{cabecera}\vspace{1cm}

\noindent D. Jesús Enrique Sierra García, profesor del departamento de Ingeniería Civil, área de lenguajes y sistemas informáticos.

\noindent Expone:

\noindent Que el alumno D. \nombre, con DNI 71705119C, ha realizado el Trabajo final de Grado en Ingeniería Informática titulado Seguimiento de líneas basado en openCV para AGVs. 

\noindent Y que dicho trabajo ha sido realizado por el alumno bajo la dirección del que suscribe, en virtud de lo cual se autoriza su presentación y defensa.

\begin{center} %\large
En Burgos, {\large \today}
\end{center}

\vfill\vfill\vfill

% Author and supervisor
%\begin{minipage}{0.45\textwidth}
%\begin{flushleft} %\large
%Vº. Bº. del Tutor:\\[2cm]
%D. nombre tutor
%\end{flushleft}
%\end{minipage}
%\hfill
%\begin{minipage}{0.45\textwidth}
%\begin{flushleft} %\large
%Vº. Bº. del co-tutor:\\[2cm]
%D. nombre co-tutor
%\end{flushleft}
%\end{minipage}
%\hfill

%\vfill

% para casos con solo un tutor comentar lo anterior
% y descomentar lo siguiente
\begin{center}
Vº. Bº. del Tutor:\\[2cm]
D. Jesús Enrique Sierra García
\end{center}

\null%\cleardoublepage
\null\cleardoublepage




\frontmatter

% Abstract en castellano
\renewcommand*\abstractname{Resumen}
\begin{abstract}
En este primer apartado se hace una \textbf{breve} presentación del tema que se aborda en el proyecto.
\end{abstract}

\renewcommand*\abstractname{Descriptores}
\begin{abstract}
Palabras separadas por comas que identifiquen el contenido del proyecto Ej: servidor web, buscador de vuelos, android \ldots
\end{abstract}

\clearpage

% Abstract en inglés
\renewcommand*\abstractname{Abstract}
\begin{abstract}
A \textbf{brief} presentation of the topic addressed in the project.
\end{abstract}

\renewcommand*\abstractname{Keywords}
\begin{abstract}
keywords separated by commas.
\end{abstract}

\clearpage

% Indices
\tableofcontents

\clearpage

\listoffigures

\clearpage

\listoftables
\clearpage

\mainmatter
\capitulo{1}{Introducción}
Los AGV, vehículos de guiado autónomo\cite{wikiagv}, actualmente son ampliamente utilizados en el mundo industrial. Su función principal es la de llevar algún tipo de carga de una zona a otra dentro de un recinto. Para ello utilizan distintos métodos de guiado, desde detección de pulsos o bandas electromagnéticas, hasta guiado óptico con cámaras\cite{guide_sistems}. 

Los métodos actuales de guiado óptico se basan en colocar una cámara muy cerca del suelo, en una posición paralela a este, de tal forma que puedan ver una "guía", una linea de algún tono que contraste con el suelo, para guiarse. El problema de esta solución, es que la imagen que la cámara es capaz de ver es un pequeño fragmento de la línea, con el que únicamente es capaz de detectar si está o no sobre la línea, permitiendo corregir la trayectoria cuando se detecta que el AGV se desvía, cuando "pierde la línea". Esta implementación es muy sensible a los cambios de luminosidad y rugosidad del suelo, haciendo que en algunos tipos de suelo industriales sea inviable.

Esta solución, se puede mejorar si aplicáramos algunos de los métodos que se están usando para el guiado de vehículos autónomos en carretera, donde no solo se detecta si el coche está dentro de los margenes del carril, sino que además son capaces de ver en perspectiva lo que tienen delante. Esto permite, además de guiar el vehículo, estimar la trayectoria que tiene que seguir en los siguientes instantes, permitiendo realizar ajustes más precisos tanto en la velocidad como en la dirección. Se podría decir que son capaces de ver su futuro, el sitio por el que tienen que pasar a continuación, por lo que son menos sensibles a los cambios de luminosidad y a la rugosidad del suelo al poder ver mucha más superficie. 

Además de esto, si el AGV puede ver lo que tiene delante, puede detectar obstáculos y objetos que se interpongan en su trayectoria, permitiendo incluso evitar colisiones.

En este proyecto, se buscará probar e implementar un software que permita realizar este tipo de guiado en perspectiva. En concreto indagaremos en la detección de líneas y el tratamiento de la perspectiva en busca de lograr ver la línea guía en dos dimensiones, para poder obtener sus dimensiones reales. 

\tablaSmall{Comparativa sistemas de guiado\cite{comp_gui}}{l c c c c c c}{comp_sist_guiado}
{ \multicolumn{1}{l}{Sistema} & F.Colocar &  F.Cambiar & Pasi. & Sucie. & Invis. & Persp.\\}{ 
Guía Magnética & Sí & Sí & Sí & Sí & Sí & No\\
Cable de inducción & No & No & No & Sí & Sí & No\\
Óptico tradicional & Sí & Sí & Sí & No & No & No\\
Óptico perspectiva & Sí & Sí & Sí & No & No & Sí\\
}

\begin{itemize}
	\item F.Colocar, facilidad de colocar.
	\item F.Cambiar, facilidad de cambiar.
	\item Pasi, pasivo, no necesita alimentación externa.
	\item Sucie, suciedad, inmune a la suciedad.
	\item Invis, invisible, se puede ver el sistema de guiado.
	\item Persp, perspectiva, podemos usarlo para ver en perspectiva.
\end{itemize}

Como podemos ver en la Tabla 1.1 el único sistema que nos puede permitir ver en perspectiva para anticipar las maniobras que ha de realizar el robot es el guiado óptico en perspectiva. 

Es cierto que tiene dos inconvenientes como es que no es inmune a la suciedad, y que la línea guía queda a la vista, el problema de la suciedad que puede llegar a ocultar la línea guía se puede solventar usando colores muy llamativos y sistemas de binarización adecuados.

Que la línea se vea  no es un gran problema, simplemente lo planteamos por que los sistemas magnéticos sí son invisibles, pero obviamente si buscamos un guiado óptico, la línea tiene que ser visible.

El resto de aspectos es igual o mejor que sus competidores, es fácil de colocar, fácil de cambiar, es pasivo(no necesita ninguna tipo de alimentación, es una linea pintada en el suelo) y nos permite ver en perspectiva.
\capitulo{2}{Objetivos del proyecto}

En este proyecto se busca desarrollar un software capaz de interpretar una imagen, distinguir en esa imagen una línea guía, y obtener la trayectoria que sigue esa guía; para posteriormente generar una serie de instrucciones capaces de dirigir a un vehículo de guiado autónomo (AGV).
Esta funcionalidad se conseguirá bien mediante el desarrollo de funciones, o empleando funciones contenidas en librerías.

También tiene el objetivo de probar en diferentes entornos todas las funciones y desarrollos que se vayan realizando, así como de documentar y registrar esas pruebas adecuadamente; con el fin último de poder reproducir el funcionamiento de una parte o de todo el software en un entorno similar al que se han realizado las pruebas, y ser capaces de obtener resultados similares a los obtenidos en esas pruebas.

\capitulo{3}{Conceptos teóricos}

\section{Implementación hardware}
En este proyecto trabajaremos con AGVs. Necesitaremos 3 elementos principales para implementarlo:

\begin{itemize}
	\item Una cámara que capte imágenes, situada adecuadamente en el AGV.
	
	\item Una ordenador que ejecute el programa.
	
	\item Un elemento para transmitir las instrucciones que genere nuestro programa al sistema de control del vehiculo.
	
\end{itemize}

Tras analizar los diferentes elementos y la forma de combinarlos, podemos obtener dos formas de organización de estos elementos, veamos ambas, y sus ventajas e inconvenienetes.
 

\subsection{Implementación total en el AGV}
En esta implementación se colocarán todos los elementos en el AGV, ejecutando el programa en el ordenador del propio vehículo.
\subsubsection{Ventajas}
\begin{itemize}

	\item No se necesita ningún elemento extra.
	
	\item La transmisión de la imagen y de las instrucciones es inmediata y por conexión física, lo que asegura una buena transmisión de la información.
	
\end{itemize} 

\subsubsection{Desventajas} 
\begin{itemize}

	\item La capacidad de procesamiento del ordenador del vehiculo no es muy grande.
	
	\item Esta implementación dificulta el control simultaneo de varios AGVs.
	
\end{itemize}

\subsection{Externalización del ordenador que ejecuta el programa} 

En esta implementación se externaliza el ordenador donde se ejecuta el programa, a un ordenador fijo, más potente.
\subsubsection{Ventajas}
\begin{itemize}
	\item El ordenador externo nos ofrece mayor capacidad de procesamiento.
	
	\item Poder ver las instrucciones y situación en tiempo real que se están mandando.
	
\end{itemize}

\subsubsection{Desventajas}
\begin{itemize}
	\item Se necesita hardware extra, además del que lleva el AGV.
	
	\item Necesitamos hardware que nos permita un buen flujo de imágenes hacia el ordenador, y una buena transmisión de las instrucciones hacia el vehículo.
	
\end{itemize} 

Para este proyecto, hemos elegido la opción de transmitir la imagen de la cámara vía inalambrica, concreamente wifi, a un ordenador externo, y ejecutar el programa en el ordenador externo. Esto nos permite tener libertad de mover la cámara a nuestro gusto, lo cual es beneficioso de cara a realizar pruebas del programa.

El alcance del proyecto no incluye el transmitir de vuelta instrucciones, ya que lo que buscamos es detectar líneas, pero esta transmisión se haría igualmente vía inalambrica, poniendo de acuerdo las instrucciones generadas por el programa y las instrucciones que interpreta el sistema controlador del vehículo.


\section{Procesado de imagen}
Para realizar una detección efectiva de la línea guía en la imagen, tenemos que realizar un procesamiento de la imagen, pasandola por distintos filtros para finalmente obtener la trayectoria que representa la línea guía.

\subsection{Binarización}
Binarizar es el proceso que nos permite distinguir ciertas partes o elementos de una imagen aplicando una serie de valores umbrales. La forma en la que se distinguen los elementos es generando una imagen binaria, donde los elementos destacados tienen un valor, y el resto de la imagen tiene otro valor.
\imagen{img_binaria}{Izquierda, imagen binaria. Derecha imagen real.} 

Primero, tenemos que entender como funciona una pantalla y una cámara, para comprender el formato de imagen más usado actualmente, RGB.

Una imagen RGB está formada por 3 matrices de valores de mismas dimensiones, una matriz para los rojos, otra para los verdes, y otra para los azules (Red-Green-Blue). Este formato, es utilizado actualmente para la representación de imágenes en pantalla, y para la captación de imágenes, ya que se ajusta al formato físico tanto de pantallas como de sensores de cámaras.
Veamos algunos conceptos relacionados con esto:

\begin{itemize}

	\item Pantalla
	
\end{itemize}

Una pantalla está formada por diminutos puntos de color, conocidos como pixeles. Cada pixel tiene 3 "luces", una para rojos, otra para verdes y otra para azules. El formato RGB da a cada una de estas luces un valor, que las hará lucir con mayor o menor intensidad.
\imagen{pixels}{Pixels de una pantalla LCD. Fuente: youtube.com}

La resolución tanto de la pantalla como de la imagen es un aspecto relevante, ya que influirá en la forma en la que la pantalla representará la imagen.
 
La resolución de la pantalla es el número de pixels de ancho, por el número de pixels de alto. La resolución de la imagen es el tamaño de la matriz de ancho, por el tamaño de la matriz de alto.

Cuando ambas son iguales, cada valor de la matriz RGB se interpretará en cada uno de los pixeles de la pantalla. En el caso de que sean diferentes, se hará un reescalado de imagen.

Los valores que pueden tomar son desde 0, apagado, hasta 255, intensidad máxima. La combinación de intensidades de las luces RGB formará los distintos colores. 
\imagen{RGB_rueda}{Combinación de colores RGB. Fuente: wikipedia.org}

Como vemos, tenemos $256^3$ combinaciones diferentes, lo que nos da un rango de 16777216 colores posibles.

\begin{itemize}

	\item Cámara
	
\end{itemize}

La forma de funcionamiento de una cámara se basa en un sensor con celdas sensibles a la luz, que permiten registrar la intensidad con la que incide la luz.
Una vez más, cada una de estas celdas se corresponderá con un una de las luces de un pixel. Para poder registrar los valores RGB se descompone la luz en sus 3 colores primarios, rojo, verde y azul. Se registra la intensidad de cada color, y se genera la imagen RGB.
\imagen{sensor_camara}{Esquema del funcionamiento del sensor de una cámara. Fuente: globalspec.com}

Una vez visto el porqué del formato RGB, podemos empezar a explicar los métodos de binarización.

\subsubsection{Binarización por luminosidad}
Este tipo de binarización consiste en generar una imagen binaria, mediante un umbral aplicado a una imagen donde podamos ver la luminosidad de cada pixel.
Para realizar este tipo de binarización necesitamos una imagen en escala de grises, donde solo tenemos luminosidad. Nuestra cámara nos devuelve imágenes RGB, por lo que debemos realizar una conversión. 

Esta conversión no es algo trivial, ya que los colores primarios de la luz, RGB, no tienen la misma luminosidad. Por lo que la solución mas simple que se nos podría ocurrir, realizar la media para cada pixel de sus valores RGB sería errónea.

En su lugar, debemos corregir la luminosidad de cada matriz de color, para igualar estas diferencias de luminosidad que por propia naturaleza los colores primarios tienen. 

Para ello, multiplicaremos cada matriz por un valor corrector.
Según la documentación de OpenCV\cite{OpenCVRGBGRAY}, la formula que nos permite hacer esta transformación es:

RGB[A] to Gray:$Y=0.299*R+0.587*G+0.114*B$

Donde Y es la imagen resultante.

Una vez obtenemos la imagen en escala de grises, podemos empezar con la binarización.

Para binarizar, primero tenemos que saber cuantos pixels hay de cada valor entre 0 y 255, para tratar de buscar grupos grandes de pixels que compartan valores similares, y así distinguir unos de otros. 

Esta información nos la da el \textbf{histograma de la imagen}.

El histograma de una imagen es la representación de la distribución de los valores de los pixels de la misma.

En el eje X tendremos de izquierda a derecha valores de 0 a 255. En el eje Y de abajo a arriba valores desde 0 hasta en máximo número de pixels encontrados en ese valor de x.

\imagen{img-histograma}{Imagen en escala de grises de la línea guía.}
\imagen{histograma}{Histograma de la Figura 3.5.}

Como podemos apreciar, hay dos picos claros en el histograma, uno correspondiente a la línea guía (pixels oscuros, pico izquierdo del histograma), y el resto correspondientes al fondo blanco (pixels claros, pico derecho del histograma), concuerda con la imágen ya que hay mayor cantidad de pixels claros que de oscuros. El umbral tiene que estar entre ambos picos para diferenciar unos pixels de otros.

La forma más simple que se nos puede ocurrir es la de escoger un valor aleatorio entre ambos picos y entonces, recorrer todos los pixels de la imagen en escala de grises, y a cada pixel asignarle un valor si esta por encima del umbral, u otro valor su esta por debajo.

Esta forma, siendo perfectamente valida, es poco recomendable si trabajamos con flujos de datos continuos, como por ejemplo vídeos, o streamings.
Normalmente en este tipo de flujos de imágenes, la luminosidad puede variar, por lo que el umbral al ser fijo, puede provocar errores en la binarización.

La solución a esto es usar un umbral dinámico, calculado en función del histograma de cada imagen.

La forma que hemos usado para calcular este umbral ha sido mediante el Algoritmo de Otsu\cite{wikiotsu}, ya implementado en OpenCV.

Lo que hace el algoritmo de Otsu es separar la imagen mediante el umbral en dos zonas. Buscamos que la dispersión dentro de cada segmento sea la mínima (que los pixels dentro de ese segmento se parezcan), pero que entre ambos segmentos sea lo máximo posible.

Para esto, inicialmente calcularemos la media aritmética de los valores de gris de toda la imagen, y después solo de cada zona del histograma. Con estos valores podemos calcular las varianzas de cada zona. 

Lo que tenemos que hacer es mantener las variazas de cada zona lo más pequeñas posibles, y conseguir que la varianza entre ambas zonas sea la máxima.

Para conseguir esto, haremos el cociente de la varianza entre las zonas y la suma de las varianzas de cada zona, buscando que este cociente sea el máximo posible. 

Veamoslo de una forma más técnica:

$K_0(t)$ y $K_1(t)$ son las zonas del histograma, separadas por el umbral t.

$p(g)$ es la probabilidad del valor gris g, donde g puede ir desde 0 hasta 255, según el formato de escala de grises que estamos usando.

La probabilidad de ocurrencia para cada zona será:
\begin{itemize}

\item En la zona que va de 0 a t:

$P_0(t)= \sum_{g=0}^{t}p(g)$ 

\item En la zona que va de t+1 hasta 255:

$P_1(t)= \sum_{g=t+1}^{255}p(g)$ o $1-P_0(t)$

\end{itemize}

Siendo $\bar{g}$ la media aritmética de los valores de gris para toda la imagen, y $\bar{g_0}$ y $\bar{g_1}$ la media para cada zona, podemos calcular las varianzas de los segmentos:

\begin{itemize}

\item En la zona que va de 0 a t:

$\sigma{_0^2} = \sum_{g=0}^{t}(g-\bar{g_0})^2p(g)$

\item En la zona que va de t+1 hasta 255:

$\sigma{_1^2} = \sum_{g=t+1}^{255}(g-\bar{g_1})^2p(g)$

\end{itemize}

Ahora necesitaremos las varianzas entre los segmentos, y la suma de la varianza de ambos segmentos.

\begin{itemize}

\item Varianza entre segmentos:

$\sigma{_zw^2} = P_0(t)*(\bar{g_0}-\bar{g})^2 + P_1(t)*(\bar{g_1}-\bar{g})^2$

\item Suma de las varianzas de cada segmento:

$\sigma{_in^2} = P_0(t)*\sigma{_0^2}(t)+P_1(t)*\sigma{_1^2}(t)$

\end{itemize}

Por último, nos queda obtener el cocientre de la varianza entre segmentos y la suma de las varianzas de cada segmento.

$Q(t) = \displaystyle\frac{\sigma{_zw^2}}{\sigma{_in^2}}$

Este es el cociente que tenemos que maximizar. El umbral será el valor de t.

En OpenCV es sencillo hacer este tipo de binarización, ya que contamos con una función que lo realiza instantáneamente, y nos devuelve tanto el umbral que ha usado, como la imagen ya binarizada.

\begin{verbatim}
umbral, img_binaria = cv.threshold(img,0,255,cv.THRESH_BINARY+cv.THRESH_OTSU)
\end{verbatim}

Indicamos que queremos hacer la binarización con el algoritmo de Otsu. El valor que se les da a los pixels por debajo del umbral es 0 y el valor por encima es 255. La imagen de origen es img.

El resultado es este:

\imagen{img_bin_otsu}{Resultado de la binarización de la Figura 3.5}

Como vemos, en esta imagen se puede diferenciar perfectamente la línea guía del resto de elementos de la imagen.

\subsubsection{Binarización por color}
Inicialmente, en el formato RGB tenemos colores, pero están separados en 3 matrices diferentes, por lo que es complicado poder binarizar de una forma efectiva. 

Existe in formato de imagen conocido como HSV\cite{hsl_hsv} (Hue, Saturation, Value), donde la matriz H corresponde con los colores puros de la imagen, la matriz S corresponde a los valores de saturación de cada pixel, y la matriz V corresponde con los valores de negro de cada pixel.

Si representamos esto en un prisma, obtenemos:

\imagen{hsv_prism}{Cilindro de colores HSV. Fuente: wikipedia.org}

En este tipo de formato tenemos la posibilidad de quedarnos con un pequeño fragmento del cilindro, dividiéndolo según el valor de H, con lo que podríamos binarizar la imagen según un rango de color concreto.

Veamos como pasar de RGB a HSV.

Lo primero que necesitamos es pasar los valores RGB a un entorno en 2 dimensiones. En RGB además de los 3 colores primarios, tenemos las mezclas de colores primarios, amarillo, magenta y cyan. Si a esto añadimos el color blanco puro, y el color negro puro, obtenemos un cubo. 
Este cubo puede ser representado en 2 dimensiones como un hexágono perfecto, usando como centro del hexágono los vértices de color blanco puro y negro puro. 
Los valores de RGB serán los que determinen la posición del pixel dentro del cubo, la cual será proyectada en dos dimensiones sobre el hexágono.

\imagen{hex_rgb_hsv}{Conversión de RGB a HSV. Fuente: wikipedia.org}

Con esto podemos obtener el ángulo del color de nuestro pixel, H. Como OpenCV trabaja con matrices de valores enteros de 8 bits, no podemos meter en la matriz valores que salgan del rango $[0,255]$ por lo que el ángulo será dividido entre 2, pudiendo estar entre 0 y 180.

El valor de negro, V, es el mayor de los valores RGB. Como OpenCV trabaja con enteros de 8 bits, este valor encaja perfectamtente.

El valor de la saturación, S, se calcula a partir del cubo RGB antes mencionado. Suponemos que el hexágono tiene un radio de 0 a 1. Al representar el color RGB en el plano, estará mas o menos alejado del radio del hexagono, este valor se conoce como Chroma, C. Para calcular la saturación, dividiremos C entre V, lo cual nos dará la saturación. Posteriormente convertiremos este valor en escala de 0 a 255, para que OpenCV pueda trabajar con ello.

Repitiendo este proceso para cada uno de los pixels de la imagen, tendremos nuestra imagen en formato HSV.

Una vez tenemos hecho esto, el proceso de binarización es sencillo. Basta con elegir en la rueda de color el ángulo de los colores que queramos. Para el ángulo 0 tenemos el rojo, 60 amarillo, 120 verde, 180 cyan, 240 azul y 300 magenta. Transformaremos el ángulo para que OpenCV lo interprete adecuadamente, simplemente es dividirlo entre 2.

Una vez tengamos el ángulo del color, estableceremos un margen de seguridad\cite{track_hsv}, para coger todos los pixels con un tono similar. Este margen ha de ser de +10 y -10 para H, y con S y V podemos jugar un poco aunque lo ideal es abarcar el máximo rango posible. Un rango recomendado es entre 100 y 255.

\imagen{img_bin_color}{Resultado de la binarización por color, escogiendo la gama de amarillos.}

En este proyecto se ha simplificado todo este proceso al máximo, ya que es algo complejo. Al usuario se le pedirá calibrar el color mediante unas ventanas de OpenCV, una donde se ve la imagen a color y otra donde se ve el resultado de la binarización. Para elegir un color simplemente tiene que pinchar encima de ese color en la imagen.

Todo esto se implementará en una clase donde se almacenará el rango de color y la funcionalidad para obtenerlo.

Finalmente, la función de OpenCV que nos permitirá realizar la binarización es:
\begin{verbatim}
inRange(matriz, umbralbajo, umbralalto)
\end{verbatim}

Donde matriz es la imagen en formato HSV, umbralbajo es el vector de valores HSV resultantes de restar a los valores del color original el margen, y umbralalto, resultante de sumar a los valores del color el margen.

\section{Calculo de la luminosidad de la imagen}
A la hora de realizar pruebas tenemos que conocer los parámetros ambientales del entorno en el que estamos, el más importante para este proyecto es la luminosidad de la imagen. Para calcular esta luminosidad hemos transformado la imagen a escala de grises, con el mismo procedimiento visto anteriormente, y hemos hecho la media de valores de la imagen.

La luminosidad máxima sería de 255, en el caso de ser una imagen donde todos los pixels sean blanco puro, y 0 en el caso de que sean negro puro. Entre estos dos valores estará la luminosidad de la imagen.


\section{Transformación de perspectiva}
Como hemos visto en secciones anteriores, la principal funcionalidad buscada en este proyecto es el poder ver la línea guía no solo debajo del AGV, también delante, en perspectiva. Para esto simplemente tenemos que colocar la cámara en una posición determinada, bastante levantada del suelo, y con un ángulo que nos permita ver la máxima perspectiva posible, sin llegar a sobrepasar el horizonte, ya que la información que necesitamos está en el suelo, la línea guía.

\subsubsection{Colocación de la cámara}

\imagen{angulo_altura}{Angulo y altura de la cámara respecto del suelo.}
\imagen{distancia}{Distancia a la que es capaz de ver la cámara en esa posición.}

Como vemos, la cámara no está a una gran altura del suelo, aún así, es capaz de ver a una distancia considerable, casi 5 veces de la altura a la que se encuentra. Por supuesto en función del ángulo está distancia será mayor o menor. 

Si tenemos en cuenta que la cámara en un AGV está colocada a unos 50cm de altura, podemos multiplicar las distancias obtenidas en la Tabla 3.1 por un factor de 6. Llegando a ver a unos dos metros y medio de distancia en el caso de colocar la cámara a 55º.

Tenemos que considerar también que al aumentar el ángulo, perdemos visión en la parte inferior de la imagen, por lo que no podremos ver exactamente lo que el vehículo tiene debajo, sino que empezaremos a ver la línea a una distancia determinada respecto de la situación real del vehículo.

\tablaSmall{Valor del ángulo y distancia a la que es capaz de ver}{l c c}{valores_angulo_distancias}
{ \multicolumn{1}{l}{Angulo} & Distancia Min & Distancia Max\\}{ 
40º & 0cm & 12cm\\
45º & 1.5cm & 19cm\\
50º & 3cm & 37cm\\
55º & 4cm & 40cm\\
60º & 5cm & Por encima del horizonte\\
}

En la Tabla 3.1 está la relación de valores. Distancia Min es la distancia entre el vehiculo y la parte inferior de la imagen. Distancia Max es la distancia hasta la que ve nuestra cámara. En el caso de colocarla con 60º, vemos que ve por encima del horizonte, por lo que no es aconsejable. Este proyecto no está preparado para este tipo de ángulo de cámara, ya que toda la información que necesitamos están el suelo, ver por encima del horizonte es perder información, y tener que realizar más filtrados de imagen, lo cual empeora el rendimiento del programa. 

\subsubsection{Obtener vista de pájaro}

El problema de situar la cámara en una posición así, es que no podemos ver las distancias reales en la imagen, puesto que lo estamos viendo todo en perspectiva. Para ver las distancias reales, tendríamos que colocar la cámara paralela al suelo y a una determinada altura, lo cual es complicado. Habría que tener una especie de brazo por delante del robot con la cámara colgando.

Por este motivo, lo que se usa para resolver este problema es una transformación de perspectiva. En concreto, una transformación que, a partir de una imagen tomada desde una situación en la que la cámara forma un cierto ángulo con el suelo, genera una imagen virtual que se vea totalmente paralela al suelo, podríamos llamarlo una "vista de pájaro".

\imagen{bird_view}{Izquierda, vista de pájaro. Derecha, vista original.}

OpenCV no tiene una función especifica para realizar esto, pero tiene una función que permite transformar perspectivas mediante unas coordenadas de origen y unas de destino\cite{persp_trans}. Las coordenadas simplemente son puntos dentro de la imagen, arrays de dos valores correspondiendo a las coordenadas X e Y.

\imagen{ej_corr_persp}{Correción mediante coordenadas}

Podemos ver como mediante las cuatro coordenadas de origen, obtenidas mediante las esquinas de la carta, se puede obtener la carta "plana", para esto simplemente tendremos que decirle al programa en las coordenadas de destino la forma que debería de tener, en este caso un rectángulo.

Una vez que conocemos como realizar está corrección mediante coordenadas, tenemos que ver como corregir la distorsión de nuestra imagen, donde se encuentra esa distorsión.

La forma más simple de llevar esto acabo es poniendo una figura geométrica ante la cámara, en nuestro caso un cuadrado. Como ya hemos visto en figuras anteriores, como por ejemplo la Figura 3.1, en la cámara no vemos un cuadrado, sino un trapecio. Esto significa que la perspectiva está agrandando los objetos contra más cerca están de la cámara, a si que para corregirla lo que debemos hacer es reducir la parte inferior de la imagen, hasta que las líneas del cuadrado sean paralelas.

Para realizar esto con la función de transformación de perspectiva de OpenCV tenemos que elegir como coordenadas de origen las cuatro esquinas de nuestra imagen, y como destino dejar fijas las esquinas superiores de la imagen, y reducir la parte inferior, moviendo las esquinas hacia el centro de la imagen\cite{coord_persp}. De esta manera corregiremos la perspectiva, pero nos quedará una imagen como la que se aprecia en la Figura 3.11, con dos triángulos negros en los extremos, que se corresponden con la reducción hecha en la parte inferior. 

\subsubsection{Relación del ángulo de la cámara con la transformación de perspectiva}

Para obtener la cantidad exacta de pixels que se han de reducir en la imagen anterior, se han llevado a cabo una serie de pruebas, con diferentes ángulos de cámara.

Se ha implementado una función que obtiene un coeficiente en función del ángulo que tiene la cámara con el suelo. Para ello utiliza la plantilla del cuadrado. Obtiene la altura del trapecio que se ve en la imagen, obtiene la diferencia entre la base y el lado superior del trapecio, y hace el cociente entre estos dos valores.

Se han anotado para cada ángulo probado la cantidad de pixels necesarios para corregir la distorsión.

Podemos ver los valores del ángulo, coeficiente y pixels a reducir en la Tabla 3.2.

Haciendo una regresión podemos obtener una función que nos relacione el coeficiente único de cada ángulo, con la cantidad de pixels a reducir.

\tablaSmall{Valores obtenidos en las pruebas de corrección de distorsión de perspectiva}{l c c}{valores_correcion_perspectiva}
{ \multicolumn{1}{l}{Angulo} & Coeficiente & Pixels\\}{ 
40º & 0.506 & 100px\\
45º & 0.530 & 105px\\
50º & 0.550 & 110px\\
55º & 0.576 & 135px\\
60º & 0.613 & 140px\\
}

Cogemos como valores de x los coeficientes, como valroes de y los pixels, y hacemos la regresión. La función resultante es:

$y = 140.3462 + (101.7432 - 140.3462)/(1 + (x/0.5605076)^64.36072)$

Esta función nos permite obtener los pixels a reducir a parir del coeficiente. El coeficiente se obtendrá a través de la función antes mencionada en la etapa de calibración del programa.



\section{Detección de bordes}
La detección de bordes es un tipo de procesado de imagen que permite obtener los bordes de los elementos de una imagen. Uno de los algoritmos más simples y más usados es el algoritmo de Canny\cite{canny_edge}. 

\subsection{Fases del algoritmo de Canny.}

El algoritmo de Canny puede ser dividido en 5 fases:
\begin{itemize}

	\item Aplicar un filtro Gaussiano para suavizar la imagen, buscando reducir el ruido de la misma. Este tipo de filtro lo que hace es desenfocar la imagen, asignando a un pixel un valor promedio respecto de los pixels que tiene alrededor. Con esto se consigue suavizar los pixels con ruido, que pueden dar lugar a errores en la detección de bordes.
	
	\item A partir de la imagen desenfocada, buscaremos los bordes. Un borde puede apuntar en muchas direcciones, el algoritmo de Canny usa cuatro filtros para detectar bordes horizontales, verticales o diagonales (0º, 45º, 90º y 135º). El operador de deteccion de bordes, devuelve la derivada en dirección horizontal y vertical, a partir de esas direcciones podemos determinar el gradiente.
	
	\item Aplicar una supresión no-máxima para deshacernos de los bordes detectados erróneamente. Para cada pixel detectado como borde se compara con los pixels en la dirección del gradiente. Si es más fuerte que los otros pixels, mantendremos el pixel, si no, lo suprimiremos.
	
	\item Tras realizar un primer filtrado por gradiente, aplicamos una binarización de doble umbral por gradiente para determinar los bordes fuertes y débiles. Se comparan los gradientes de los bordes resultantes, etiquetándolos como bordes débiles, si están entre los dos umbrales, borde fuerte si están por encima del umbral mayor, y desechando los que estén por debajo del umbral menor.
	
	\item Encontrar los bordes por histéresis, buscando  la cercanía con algún otro borde fuerte. Suprimir los débiles no conectados a bordes fuertes.
	
\end{itemize}

\subsection{Implementación de Canny en OpenCV.}

En OpenCV esta está implementada como\cite{canny_doc}

\begin{verbatim}
bordes = cv2.Canny(img, th1, th2, aperture, L2gradient)
\end{verbatim}

\begin{itemize}

	\item img, imagen sobre la que se va a aplicar el algoritmo.
	
	\item th1 primer valor para reaizar la histéresis.
	
	\item th2 segundo valor para realizar la histéresis. 
	
	\item aperture, tamaño de la apertira para el operador Sobel, usado para destacar los bordes de la imagen.
	
	\item L2gradient, bandera booleana, en verdadero aplica la regla L2 más precisa para calcular la magnitud del gradiente de la imagen. Si está a falso aplica la regla L1.
	
\end{itemize}

\imagen{edge_detection}{Detección de bordes a partir de la imagen binaria.}

\section{Generación de la trayectoria}
Una vez obtenidos los bordes de la línea, obtener la trayectoria es algo sencillo. Basta con calcular el punto medio entre los dos bordes, y asignarle un valor de 255 para registrarlo como línea. Se ha implementado una función que recoge esta funcionalidad. 

\imagen{tray_bordes}{Detección de la trayectoria a partir de los bordes.}

\section{Sistema de guiado básico}
El problema del cual surgía este proyecto, es que los sistemas actuales de guiado solo detectan si nos salimos o no de la línea guía. Teniendo la trayectoria ya definida, podemos crear algo con un funcionamiento parecido, pero a partir de la trayectoria. 

\imagen{guiado}{Pantalla del sistema de guiado básico.}

\subsection{Posición del vehículo respecto de la línea}
Lo que haremos será detectar la posición de la trayectoria en la zona inferior de la imagen, y ver si está a la izquierda o a la derecha del centro de la imagen. La trayectoria es la línea azul, el centro de la imagen esta marcado con el indicador verde en la zona inferior de la imagen.

Con estos dos elementos ya podemos tener un sistema muy básico que nos diga si estamos a la derecha o a la izquierda de la trayectoria ideal.

Partiendo de esta idea, nos bastaría con generar una instrucción que enviaríamos al controlador del vehículo, diciéndole hacia donde se tiene que dirigir, si estamos dejando la línea a la izquierda del centro, le mandaríamos girar a la izquierda, si fuese a la derecha, le mandaríamos girar a la derecha.

\subsection{Mejoras de en generación de instrucciones de giro}

\subsubsection{Mejora en el ángulo de giro}
Una gran mejora que podemos realizar, si está soportada por el controlador del vehículo, es poder controlar el ángulo de giro en función de lo que nos estemos desviando del centro. 

Es irreal que ante una mínima desviación, y una gran desviación, la orden que le demos al vehículo sea igual, ya que en el primer caso con una leve corrección sería suficiente, mientras que en el segundo necesitaremos una corrección más grande.

La solución planteada a este problema es usar una función que nos calcule el ángulo de giro en función de la cantidad de pixels que nos estemos alejando del centro de la imagen. Mediante regresión, podemos obtener está función, dando como valores de X la cantidad de pixels de desviación y como Y el ángulo deseado. A menor desviación, menor ángulo, si estamos sobre el centro ideal, ángulo de giro 0.

\subsubsection{Mejora en la frecuencia de corrección}
Otro problema que nos podemos plantear es si de verdad necesitamos estar corrigiendo continuamente al vehículo, o podemos hacer un guiado más natural.

Cuando las personas conducimos, no estamos todo el rato buscando el centro del carril, ya que sería casi imposible, lo que hacemos es tener unos margenes "seguros" donde nos podemos colocar, y en el momento en el que nos acercamos al limite de uno de esos margenes es cuando corregimos la trayectoria.

La idea es llevar esto al sistema de guiado, establecer unos márgenes a la derecha y a la izquierda del centro ideal, y mientras la trayectoria esté dentro de estos márgenes, mandar al vehículo seguir recto. En la Figura 3.17 podemos ver estos márgenes, las líneas rojas en la parte inferior de la imagen, próximas al centro.

Los márgenes se generarán en la etapa de calibración del programa, se solicitará al usuario un valor entre 0 y 0.3, y se multiplicará este valor por la resolución horizontal de la imagen. Esto nos dará un valor en pixels, y estos pixels serán la distancia del margen a la izquierda y derecha del centro. Obviamente si damos de margen 0, no tendremos margen y estaremos en la situación en la que el programa nos corrige continuamente. El limite superior es 0.3 ya que si superamos este valor, el margen sería demasiado grande, llevan al programa hasta casi perder la línea.

La funcionalidad concreta es, mientras la línea de la trayectoria este entre estos márgenes, el programa seguirá generando la instrucción para seguir recto. En el momento que nos salgamos de uno de estos márgenes, el margen pasará a ser el centro, por lo que el programa nos corregirá hasta volver al centro. Una vez que lleguemos al centro, el margen volverá a su posición anterior.

En este proyecto las instrucciones que genera el sistema de guiado se muestran por pantalla. Podemos ver un indicador númerico con el ángulo de giro y un indicador textual que puede tomar los valores:

\begin{itemize}

	\item RI, recto izquierda, indica que el ángulo se tiene que interpretar como un giro a la izquierda. Para ángulos menores o iguales de 45º.
	
	\item RD, recto derecha, indica que el ángulo se tiene que interpretar como un giro a la derecha. Para ángulos menores o iguales de 45º.
	\item I, similar a RI pero para ángulos mayores de 45º.
	
	\item D, similar a RD pero para ángulos mayores de 45º.
	
\end{itemize}

Para llevar esto a una implementación real, tendríamos que ponernos de acuerdo con el sistema controlador del vehículo en la estructura de las instrucciones generadas, por ejemplo, un array de valores, una tupla... 
En este proyecto como no tenemos conexión real con un vehículo, las instrucciones se generan para ser visualizadas en pantalla.

\imagen{guiado_fin}{Comparación de situaciones en el sistema de guiado.}

En la Figura 3.18 podemos ver como en el momento en el que la línea de trayectoria sobrepasa uno de los márgenes, en concreto el izquierdo, el margen se mueve hacia el centro, y el programa comienza a generar instrucciones para corregir la trayectoria, en este caso nos manda realizar un giro hacia a la izquierda. El ángulo ira disminuyendo según el centro de la imagen se vuelva a acercar a la trayectoria, volviendo a la situación donde genera la instrucción de ir recto.


\section{Sistema de guiado virtual}
Este sistema de guiado sería el siguiente avance. Consistiría en recrear virtualmente la situación en la que está el vehículo. 

Con el sistema de guiado básico, ademas de podernos guiar de una forma simple, no perdemos la trayectoria de vista. Lo siguiente que tenemos que hacer es aplicar algún tipo de transformada a la imagen binaria de la trayectoria, para obtener el polinomio aproximado que represente esa trayectoria. Una vez hecho esto, podemos generar un entorno virtual en dos dimensiones según la posición en la que se encuentra el vehículo y la trayectoria que debe seguir.

A partir de este entorno virtual, podemos no solo guiarnos sino preparar el vehículo para la trayectoria que tiene que realizar a continuación, ajustar la velocidad si se aproxima una curva muy pronunciada, ajustar el ángulo de giro para igualar el radio de giro de la curva...

Si volvemos un poco atrás, concretamente a la Tabla 3.1, podemos obtener las medidas que necesitaríamos para representar de forma virtual la situación del vehículo. Veamos un ejemplo de este entorno virtual:

\imagen{guiado_virtual}{Entorno virtual para el sistema de guiado virtual.}

En esta imagen podemos ver como las medidas de la tabla nos pueden ayudar a saber la distancia que hay entre el vehículo (Punto rojo) y el comienzo de la trayectoria(Distancia Min), y la distancia máxima hasta la que podemos ver(Distancia max). La trayectoria está representada mediante la línea curva azul.

Esta sería la representación virtual del entorno, y nos ayudaría a generar unas instrucciones de guiado mucho más precisas. 

Este sistema de guiado se escapa del alcance de este proyecto, donde buscamos el llegar a tener los elementos necesarios para recrear este entorno virtual. La trayectoria en vista de pájaro con sus dimensiones reales, la situación del robot y un sistema para mantener la trayectoria lo más centrada posible.

\capitulo{4}{Técnicas y herramientas}

\section{Técnicas de desarrollo}

\subsection{Metodología ágil}

Este proyecto se ha resuelto mediante metodología ágil. 

Cabe destacar que este proyecto no es un proyecto software común en el que tenemos unos requisitos iniciales mas o menos claros, y tenemos alguna idea del final. En este caso lo que buscamos es probar diferentes códigos, herramientas, librerías para satisfacer un primer requisito esencial, la detección de lineas en el suelo. Si conseguimos satisfacer ese requisito, podremos proponer nuevos requisitos en base al resultado obtenido.

Entonces, no hay un objetivo final claro, si bien, lo ideal seria avanzar lo máximo posible, también se puede dar el caso de probar y probar diferentes elementos, y no llegar a obtener un resultado esperado.

Por ser un proyecto software, las ventajas de usar una metodología ágil son claras:
\begin{itemize}
	\item Las personas que participan en el desarrollo adquieren roles diferenciados.
	\item El desarrollo es incremental, por etapas. Al final de cada etapa se entrega una parte del software funcional al cliente.
	\item Partimos de unos requisitos básicos y podemos ir añadiendo más a lo largo del desarrollo, o modificar alguno de los existentes.
	\item La comunicación con el cliente se hace a lo largo de todo el desarrollo, favoreciendo que el producto final sea lo que esperaba.
\end{itemize}

Una de las metodologías ágiles más usadas es Scrum, cuyo principio clave es que los clientes pueden cambiar de idea sobre lo que quieren y necesitan. 
En Scrum hay 3 roles diferenciados:
\begin{itemize}
	\item Product Owner: es el cliente. Propietario de la idea o proyecto a desarrollar. 
	\item ScrumMaster: experto en Scrum, ayuda tanto al Product Owner como al Equipo Scrum a alcanzar los objetivos finales del proyecto.
	\item Equipo Scrum: equipo de personas que va a realizar el proyecto. Es conveniente que sea un equipo multidisciplinario, en el que cada persona sea experta o conocedora de un campo diferente a las demás. El equipo debe de ser de 5 a 9 personas, para favorecer la comunicación y fluidez del desarrollo.
\end{itemize}

Se basa en 3 principios básicos, en comparación con metodologías antiguas:
\begin{itemize}
	\item Favorecer la comunicación del equipo en lugar de la excesiva documentación.
	\item Buscar la calidad del resultado en los conocimientos de las personas que han participado en el desarrollo, y no en los procesos empleados para hacerlo.
	\item Solapar las fases de desarrollo en lugar de hacerlas de forma secuencial o en cascada.
\end{itemize}

Para este proyecto en concreto, el alumno tomara el rol de Equipo Scrum y ScrumManager. El profesor actuará como cliente.

Como antes se ha explicado, este no es un proyecto software común, por lo que partiremos desde un requisito inicial, y cuando el cliente lo consideré valido y completado, propondrá mas requisitos.

En cada sprint se explorarán diferentes herramientas en busca de un objetivo, se harán pruebas con alguna de esas herramientas, estableciendo los parámetros iniciales de cada prueba, el resultado esperado y obtenido. El software evolucionará con cada prueba.

\section{Herramientas de documentación}

\subsection{LaTeX}

Latex es un sistema de composición de texto, que busca la creación de documentos con una alta calidad tipográfica.

Para la documentacion de este proyecto se ha usado TexMaker, que es gratuito bajo la licencia GPL.

Página oficial editor Latex: http://www.xm1math.net/texmaker/

También se ha usado MikTex como implementación de Latex para windows. En su página dicen que por ser un conjunto de paquetes, no tienen una licencia concreta como en otros casos, aun así siguen las directrices de la FSF, y establecen una serie de pautas de uso distribuición y modificación.

Página oficial MikTex: https://miktex.org/

\section{Herramientas de gestión}

\subsection{Trello}
Trello es una herramienta online de gestión de tareas intuitiva y simple. Podemos crear diferentes tableros y crear tareas dentro de estos. La vamos a usar principalmente para comunicarnos entre profesor y alumno. Crearemos 3 tableros, tareas por hacer, tareas en desarrollo y tareas completadas.

Página de trello: https://trello.com/

\subsection{GitHub}
Github es una plataforma online para albergar proyectos desarrollados mediante un sistema de control de versiones git. 
Es una herramienta mundialmente conocida en el mundo del desarrollo de software, y una de las más usadas para este tipo de propósitos.

Principalmente tiene 3 tipos de elementos:
\begin{itemize}
	\item Milestones: Son puntos temporales de importancia dentro de la planificación del proyecto.
	\item Issues: Tareas a realizar. Pueden ser asignadas a um Milestone. 
 	\item Commits: Son el elemento que permite ir realizando cambios en el repositorio: añadir elementos, modificar, eliminar... Sirven para ir resolviendo las tareas que vayan surgiendo.
\end{itemize}

La usaremos para albergar el código del proyecto, y para aplicar metodologías ágiles y scrum. Los Sprints los representaremos con Milestones. Las tareas que vayan surgiendo en cada Sprint como Issues, y la resolución de las diferentes tareas se harán con commits. He elegido esta herramienta por haberla usado en alguna asignatura en la carrera y estar familiarizado con ella, y por ser uno de los repositorios de git más usados.

\subsection{ZenHub}
Es una extension para Google Chrome. Extiende las posibilidades de GitHub, añadiendo elementos como tableros para organizar las tareas, nuevas medidas para evaluar el desarrollo del proyecto y gráficas asociadas a estas medidas. He elegido esta herramienta por ser un gran añadido a GitHub, y por haberla usado en alguna asignatura de la carrera.


\section{Herramientas de desarrollo}

\subsection{Python}
El lenguaje de programación usado para este proyecto es python. Ha sido elegido, primero porque las librerias de OpenCV están para este lenguaje, lo cual es un requisito imprescindible; y segundo por la facilidad de uso, la versatilidad y la gran cantidad de librerias externas con las que cuenta.

Concretamente se ha usado python en su versión 2.7.5.
Nos hubiera gustado dar el salto a python 3, pero no hay librerías de OpenCV de forma oficial para esta versión, solo librerías traducidas de python 2 a python 3, por usuarios, por lo que no contamos con un soporte oficial y nos podemos encontrar con errores y sorpresas desagradables.


\subsection{PyLint}
PyLint es una herramienta para analizar y encontrar errores en códigos python. Se puede añadir a IDEs de forma sencilla, y se pueden personalizar los distintos avisos que proporciona. 
Además permite refactorizar código y cuenta con los estándares para el desarrollo de códigos en python.

\subsection{OpenCV}
OpenCV es una librería software de vision aritificial y machine learning. Nos proporcionará todos los elementos software necesarios para poder obtener y procesar imágenes. 
OpenCV se distribuye bajo licencia BSD. 
Esta librería está implementada para varios lenguajes de programación, aunque la elección fue python por la facilidad de uso, la ligereza del lenguaje, y las posibilidades de extension que ofrece.

\subsection{Visual Studio Code}
Visual Studio Coce es el IDE elegido para este proyecto. Es un IDE sencillo, compatible con una grán variedad de lenguajes de programación y compatible con Linux, Windows y MacOS.
 
Se pueden añadir extensiones de una forma muy sencilla dentro del propio IDE, de hecho es el propio IDE el que nos sugiere instalar algunas extensiones dependiendo del lenguaje de programación que estemos usando. 

En mi caso, en cuanto detecto que en mi repositorio había ficheros python, me recomendó instalar la extension de python, la cual incluye pylint, el debugger de python y alguna otra herramienta más.

Además de todo esto, es compatible con git de una forma muy sencilla, simplemente abrimos la carpeta de nuestro repositorio y ya entra en funcionamiento el control de versiones.

Por todo esto, considero que esta herramienta es la más adecuada y la que mejor se adapta a mis requisitos para realizar este proyecto.

Su propietario es Microsoft y sus licencias son:
\begin{itemize}
	\item Codigo fuente, licencia del MIT.
	\item Binarios, licencia FreeWare.
\end{itemize}

\subsection{IPWebcam}
Aplicación android, con dos versiones, gratuita y de pago. Actualmente se usa la versión gratuita.
Su funcionalidad es crear un streaming online en directo, accesible desde una url. Es útil puesto que me permite tener la cámara totalmente libre, sin ataduras, pudiéndola mover como desee.

A la hora de realizar pruebas esto es una gran ventaja, en lugar de estar moviendo el ordenador portátil de un lado a otro con la webcam propia del equipo, o con una cámara conectada por cable; simplemente puedo tener el móvil en un pequeño soporte con ruedas y moverlo libremente a lo largo de la linea. Además me permite tenerlo en diferentes entornos con diferentes luminosidades. Y también hay que tener en cuenta que la cámara de un smartphone es mucho mejor que la de un ordenador portátil, en cuanto a calidad de imagen y captación de luz, por lo que ofrece mayor versatilidad de cara a trabajar en diferentes entornos. Es por esto que he elegido un dispositivo como un smartphone y esta aplicación para la obtención de la imagen con la que voy a trabajar en el proyecto.





 

\capitulo{5}{Aspectos relevantes del desarrollo del proyecto}

\section{Definición como proyecto de investigación}

Lo primero a destacar es que este proyecto no es un proyecto software común en el que tenemos unos requisitos iniciales mas o menos claros, y tenemos alguna idea del producto final. En este caso lo que buscamos es probar diferentes códigos, herramientas, librerías para satisfacer el requisito principal, la detección de lineas en el suelo. Si conseguimos satisfacer ese requisito, podremos proponer nuevos requisitos en base al resultado obtenido. En definitiva, es más un proyecto enfocado en investigación, pero sin olvidarnos de que tenemos que crear un producto funcional que recoja los resultados de las investigaciones realizadas.

\section{Metodología ágil}
Como antes hemos explicado, este proyecto tiene una gran parte enfocada a la investigación, para adaptar esto a las metodologías ágiles, hemos tenido que:

\begin{itemize}
	\item Modificar la duración de los sprints hasta dar por satisfactoria la investigación realizada en el sprint. 
	\item En base a los resultados de cada investigación, proponer nuevos requisitos, sin tener un objetivo final.
\end{itemize}

Por supuesto, el resto de líneas de trabajo de las metodologías ágiles se han seguido cumpliendo:

\begin{itemize}
	\item Reuniones al finalizar cada sprint.
	\item Organizar tareas en un tablero canvas.
	\item Entregar y realizar una demostración de la investigación al finalizar cada sprint.
\end{itemize}

\section{Relación con conocimientos obtenidos en la carrera}

\subsection{Conocimientos en Python}
Lo primero a destacar es el uso de Python, lenguaje de programación visto en varias asignaturas de la carrera, como Algoritmia, Sistemas Inteligentes o Minería de Datos. 

En estas asignaturas también se han usado algunas de las librerías empleadas en este proyecto, como MatPlotLib o NumPy.


\subsection{Conocimientos de procesamiento de imágenes}
Los conocimientos de algoritmos de visión artificial, se han obtenido en la asignatura de Hardware de Aplicación Específica. Aunque en esta asignatura no se usa OpenCV, sino Matlab, la gran mayoría de funciones que se usan en Matlab para realizar el procesamiento de imagen, tienen una replica exacta en OpenCV.

\subsection{Idea del diario de investigación}
La idea del diario de investigación ha surgido de la asignatura de Fundamentos Físicos de la Informática, en la que usábamos un cuadernillo donde íbamos documentando las diferentes prácticas y los resultados obtenidos. Se puede ver aquí\cite{diario}.

\section{Conocimientos externos a la carrera}

El problema de la transformación de perspectiva es complejo, y no se ha visto durante la carrera, por lo que he tenido que buscar información por internet. La información técnica de este tipo de problemas se sale de los conocimientos propios de un informático, están más destinados a personas especializadas en Óptica, aún así, hay tutoriales y otros proyectos donde se trata este tipo de problema desde un punto más amigable para un informático. 

Un tutorial que me ayudo a comprender como realizar la transformación de perspectiva es este\cite{coord_persp}. Aunque en ese tutorial colocan la cámara viendo por encima del horizonte y tienen que realizar unos ajustes extra para obtener la vista solo desde el horizonte hacia abajo. Aun así tiene conceptos como el de reducir la parte inferior de la imagen, que han sido usados en este proyecto.



\section{Restricciones del proyecto}

\tablaSmall{Resolucion 800x600 y Fps}{l c c c}{resolucion_fps800600}
{ \multicolumn{1}{l}{Calidad streaming} & Min Fps &  Max Fps & Media Fps\\}{ 
100 & 1.43 & 6.85 & 3.82\\
75 & 4.67 & 16.13 & 11.49\\
50 & 2.15 & 19.23 & 12.04\\
}

\tablaSmall{Resolucion 640x480 y Fps}{l c c c}{resolucion_fps640480}
{ \multicolumn{1}{l}{Calidad streaming} & Min Fps &  Max Fps & Media Fps\\}{ 
100 & 2.28 & 7.3 & 5.57\\
75 & 1.07 & 23.26 & 15.73\\
50 & 2.47 & 22.73 & 16.13\\
}

\tablaSmall{Resolucion 176x144 y Fps}{l c c c}{resolucion_fps176144}
{ \multicolumn{1}{l}{Calidad streaming} & Min Fps &  Max Fps & Media Fps\\}{ 
20 & 5.26 & 200.00 & 58.20\\
}

No tenemos que olvidar que el fin último de este proyecto es implementarlo en un entorno industrial, donde todo tiene que estar meticulosamente controlado. 

Un requisito importante que se impuso desde el principio del proyecto es que un AGV necesita recibir una instrucción cada 2 centésimas de segundo, realizando una operación simple, obtenemos que necesitamos procesar 50 imágenes por segundo para poder satisfacer este requisito.

El hardware con el que cuento no es demasiado potente, ni esta específicamente diseñado para cumplir este requisito, como si lo estaría en una implementación real. Aun así, al principio del proyecto se realizaron unas pruebas de resolución y calidad de imagen, en busca de esta tasa de fotogramas.

Los resultados se muestran en las tablas: Tabla 5.4, Tabla 5.5 y Tabla 5.6.

Una vez obtenida la tasa de frames deseada, se empezó con la siguiente fase del proyecto.

\section{Fases del proyecto}
Las fases de este proyecto partían de un requisito a cumplir, realizando un trabajo de investigación, y posteriormente generando una implementación y realizando pruebas en el entorno. 

\subsection{Binarización}
\subsubsection{Binarizacion por luminosidad}
Inicialmente se planteó la binarización por luminosidad, se investigaron diferentes algoritmos como Otsu, Umbral Adaptativo por Vecindad y Umbral fijo, siendo el más efectivo el algoritmo de Otsu. 

El algoritmo de umbral fijo no era capaz de aguantar la variación de luminosidad. 

El algoritmo de Umbral Adaptativo por Vecindad no está preparado para este tipo de problemas, está enfocado a usarse en imágenes donde la luz incide de forma desigual, siendo capaz de calcular un umbral en función de la luminosidad de los vecinos del pixel que se este procesando.

El algoritmo de Otsu calcula el umbral en función de la luminosidad de la imagen, lo cual es ideal para este proyecto ya que se puede adaptar a las variaciones de iluminación. En la versión final del proyecto, se muestra tanto la luminosidad de la imagen, como el umbral de Otsu que se está usando para binarizar la imagen.

\imagen{lum_umbral}{Pantalla donde se muestra tanto la luminisdad de la imagen como el umbral de binarizacion generado por el algoritmo de Otsu.}

Bajo la experiencia de todas las pruebas realizadas, la luminosidad medida de está forma puede llegar a bajar hasta a valores de 2, donde prácticamente toda la imagen es negra, y llegar hasta valores de 200, donde la imagen es blanca entera. Llegar a generar una imagen en un entorno real donde la cámara capte el color blanco puro es prácticamente imposible, por esto nunca llegaremos a acercarnos a valores de 255. 

\subsubsection{Binarizacion por color}

En etapas más avanzadas del proyecto se probó la binarización por color, mediante el formato de imagen HSV obteniendo buenos resultados, y una implementación más realista, de cara a que no es necesario que el fondo donde se coloca la imagen sea uniforme, sino que puede ser de cualquier tipo, siempre que no sea del mismo color o similar que la linea que intentamos detectar. Aun así este tipo de binarización es mucho más sensible a las variaciones de luz, por lo que también tiene esa desventaja.

\subsection{Transofmación de perspectiva}
Para realizar esto primero se implementaron unas funciones que tenían cierta lógica, pero no acabaron funcionando como se esperaba, aunque sirvieron como punto de partida para usar algunas de las otras funciones que OpenCV tiene para resolver este problema.

\subsubsection{Función inicial}

Inicialmente se creó una función que solo trabajaba con imágenes binarizadas. Lo que hacía era mediante la plantilla del cuadrado, calculaba un "coeficiente de reducción" que servía para saber la cantidad de pixels a reducir en cada fila de la imagen. 

Lo que hacemos es a partir del trapecio que se genera al ver el cuadrado en perspectiva, calcular la diferencia entre su base y su lado superior. Esta diferencia es la que tenemos que "solventar" para hacer que sean iguales, y por tanto un cuadrado. 
El trapecio está formado por filas de pixels, por lo que podemos dividir esa diferencia en pixels de los lados entre el número de filas que tiene de alto el trapecio, lo que nos dará la cantidad de pixels que tenemos que reducir en cada fila.

Por ejemplo en la fila donde está el lado menor del trapecio, la fila 0, tendremos que reducir 0 * coeficiente de reduccion, 0 pixels.

Para la fila de la base, tendremos que reducir el número de fila donde se encuentre (altura del trapecio) * coeficiente, lo que nos dará exactamente la diferencia entre los dos lados antes calculada, dejando del mismo tamaño el lado superior e inferior.

Las filas del medio del trapecio seguirán el mismo algoritmo, dejando a todas iguales al lado menor del trapecio, y por tanto dejandolo en forma cuadrada.

Este algoritmo a pesar de que funciona, tiene un problema y es que no guarda la relación de aspecto de los objetos del mundo real. El cuadrado que debería ser cuadrado, es rectangular, a si que no nos vale.

Pero lo que si hemos conseguido es un algoritmo capaz de generar un coeficiente único para cada ángulo que forma la cámara con el suelo.

\subsubsection{Funcion de transformación de perspectiva de OpenCV}

Una vez podemos distinguir el ángulo con el que estamos, podemos relacionar este ángulo con las funciones de OpenCV, pudiendo generar una transformación de perspectiva única y adecuada para cada ángulo.

\subsection{Obtener la trayectoria}
Esta fase es la más inmediata, solo necesitamos los bordes de la línea con la distorsión por perspectiva resuelta, y obtener el punto medio entre los bordes.

Una vez generada la trayectoria, se pueden aplicar diferentes métodos para estimar el polinomio capaz de adaptarse a la forma que tenga la trayectoria.


\subsection{Sistema de guiado}
Aunque lo ideal es recrear virtualmente el entorno donde está el vehículo, como hemos visto en secciones anteriores, esta funcionalidad es realmente compleja y se escapa del ámbito de este proyecto. Aun así, se puede generar un sistema de guiado básico, con la parte más cercana de la trayectoria que tiene el vehículo. 

La idea inicial era hacer un sistema que simplemente nos dijera si estamos a la izquierda o a la derecha de la trayectoria, pero finalmente se añadieron las funcionalidades del calculo del ángulo de giro en función de la desviación, y el margen de seguridad donde permitimos que el robot siga recto aunque no este exactamente sobre la trayectoria, corrigiendole cuando se salga de este margen.

\section{Pruebas}
En cada fase se han realizado pruebas y se han dejado expuestas en vídeos en el repositorio.

Las pruebas están organizadas según la investigación:

\begin{itemize}
	\item Una prueba inicial con el objetivo de obtener 50 Fps. Podemos verlas en los siguientes vídeos:
	
	\href{https://youtu.be/r73s9I_8SiI}{Resolución 800x600 calidad 50}
	
	\href{https://youtu.be/rdv0COnl-G4}{Resolución 800x600 calidad 75}
	
	\href{https://youtu.be/tGip7sWjD8E}{Resolución 800x600 calidad 100}
	
	\href{https://youtu.be/saFDDQNjsEs}{Resolución 640x480 calidad 50}
	
	\href{https://youtu.be/a3ZxLQB37k0}{Resolución 640x480 calidad 75}
	
	\href{https://youtu.be/zu8rWr31fuQ}{Resolución 640x480 calidad 100}
	
	\href{https://youtu.be/WzJjXKeKpaI}{Resolución 176x144 calidad 20}	
	
	Lo que se ha ido haciendo es modificar la resolución y calidad de la imagen hasta que hemos obtenido la tasa de frames deseada, 50 fps. Esta tasa se ha alcanzado en la prueba de resolución 176x144, calidad del streaming 20.
	
	\item Pruebas de luminosidad, en busca de implementar una función que calcule la luminosidad de la imagen. Podemos verlas en el siguiente vídeo:
	
	\href{https://youtu.be/CuTADdioU4g}{Prueba de la función que calcula la luminosidad}	
	
	En este caso hemos modificado la luminosidad del ambiente en el que se encontraba la cámara para ver como la función se comportaba. El resultado fue satisfactorio, ya que vemos que cuando disminuimos la luminosidad, la funcion lo detecta adecuadamente.
	
	\item Pruebas de binarización, tanto por luminosidad como por color. Podemos verlas en los siguientes vídeos:
	
	\href{https://youtu.be/Gu_HvPr3cCU}{Binarizar con umbral fijo}
	
	\href{https://youtu.be/NCWQ2I-1J_o}{Binarizar con umbral adaptativo de vecindad Media}
	
	\href{https://youtu.be/XPxqk3kNXKY}{Binarizar con umbral adaptativo de vecindad Gausiano}
	
	\href{https://youtu.be/4CW4r9qdo4Y}{Binarizar por luminosidad con Otsu}
	
	\href{https://youtu.be/Rks5PkxDeZQ}{Binarizar por color}
	
	La funcionalidad buscada en esta prueba es la de obtener una función que se comportara bien a la hora de binarizar. Hemos encontrado dos funciones capaces de esto, binarización mediante el algoritmo de Otsu, poco sensible a las variaciones de luminosidad, binarización por color, sensible a las variaciones de luminosidad. En condiciones de luminosidad constante, ambas funciones se comportan bien.
	
	\item Prueba de calculo del coefciente para resolver la distorsión de perspectiva. Podemos verla en el siguiente vídeo: 
	
	\href{https://youtu.be/Y9QgFIMXSiU}{Prueba de calculo del coeficiente}
	
	En esta prueba lo que buscamos es realizar un test y anotaciones de los diferentes coeficientes que devuelve la función que calcula el coeficiente para relacionarlo con el ángulo que tiene la cámara respecto del suelo.	
	
	\item Pruebas de la resolución de la distorsión de perspectiva. Podemos verlas en los siguientes vídeos:
	
	\href{https://youtu.be/rUQrHbA_H8A}{Tomar medidas del ancho de la línea}
	
	\href{https://youtu.be/ClA-Ei8f18o}{Comprobar medidas de la linea}
	
	\href{https://youtu.be/m7t4dcfuFeo}{Prueba a medir el ancho variando la luminosidad}
	
	\href{https://youtu.be/Y9QgFIMXSiU}{Correción de la distorsión de perspectiva}
	
	\href{https://youtu.be/Ae5aTr3mg6k}{Encontrar la trayectoria con la distorsión resuelta}
	
	\href{https://youtu.be/wc00piNmfyI}{Comprobar que la imagen corregida guarda las proporciones de la realidad midiendo ángulos 1.}
	
	\href{https://youtu.be/T0djIwql6G8}{Comprobar que la imagen corregida guarda las proporciones de la realidad midiendo ángulos 2.}
	
	\href{https://youtu.be/pRE8zerfp7E}{Ver el coeficiente calculado con el ángulo}
	
	En esta serie de pruebas se busca realizar la asociación del coeficiente calculado con la transformación de la perspectiva. El resultado es satisfactorio ya que conseguimos implementar esta funcionalidad. 
	
	Además se realizan pruebas para comprobar que la imagen con la perspectiva en vista de pajaro guarda las mismas relaciones de tamaño que en la vida real. Para ello se ha creado una plantilla con ángulos, se ha colocado delante de la cámara, y se han medido esos mismos ángulos en la imagen en vista de pájaro, comprobando que son exactamente el mismo ángulo que en la realidad. El resultado es satisfactorio.
	
	\item Pruebas del sistema de guiado. Podemos verlas en los siguientes vídeos:
	
	\href{https://youtu.be/ODr3q91MLk0}{Prueba del sistema de guiado}
	
	\href{https://youtu.be/8j69cDydXfA}{Sistema de guiado final con todos los elementos visualizables}
	
	En estas pruebas vemos que el sistema de guiado se comporta de la forma esperada, por lo que son satisfactorias.
	
\end{itemize}

La función de luminosidad se ha usado en todas las etapas de procesado de la imagen, para saber las condiciones del entorno y poder reproducir la prueba bajo las mismas condiciones.

\capitulo{6}{Trabajos relacionados}

Este apartado sería parecido a un estado del arte de una tesis o tesina. En un trabajo final grado no parece obligada su presencia, aunque se puede dejar a juicio del tutor el incluir un pequeño resumen comentado de los trabajos y proyectos ya realizados en el campo del proyecto en curso. 

\capitulo{7}{Conclusiones y Líneas de trabajo futuras}
A continuación se exponen las conclusiones obtenidas en la realización del trabajo, y las posibles lineas de trabajo futuras. 

\section{Conclusiones}
Para finalizar está memoria, veremos algunas de las conclusiones obtenidas al realizar este proyecto:

\begin{itemize}
	\item Lo primero, me gustaría dar énfasis a algo que parece trivial, pero en realidad es la parte del proyecto que más conflictos puede generar. El trabajar con flujos de datos continuos y en directo. No solo tenemos que transmitir el vídeo, tiene que llegar con una latencia mínima, con una calidad que permita trabajar con ello, a unos Fps determinados, y sin parones repentinos. Además, podemos enfrentarnos a entornos no controlados que generen dificultad extra en las diferentes fases del proyecto. Esta es una de las conclusiones más relevantes, ya que es algo que inicialmente no te planteas, pero según avanzas en el proyecto te das cuenta de la importancia de ello.
	
	\item A continuación destacar la dificultad para alcanzar una robustez lo suficientemente grande como para poder usar este tipo de algoritmos y programas en un entorno profesional. El proceso de binarización es realmente complejo, pudiendo usar gran cantidad de filtros en diferentes ordenes para llegar a obtener la imagen deseada. Eso sí, siempre buscando la mayor optimización posible, ya que no debemos olvidarnos de que este proyecto tiene que funcionar a tiempo real y con una tasa de frames determinada.
	
	\item Por último, decir que el enfrentarme a un proyecto no solo de desarrollo software, sino también de investigación ha sido complicado. Aun así, los beneficios superan al esfuerzo realizado, ya que he aprendido no solo a desarrollar un proyecto software, sino a investigar en los campos que abarca este proyecto, investigar proyectos similares, documentar una serie de pruebas en base a realizar la investigación y también llevar un diario de investigación en el que ir anotando los progresos realizados.
	
\end{itemize}

En definitiva, la experiencia ha sido positiva, y muy enriquecedora para mi futuro como desarrollador de software, dándome también conocimientos sobre como llevar a cabo una investigación.
 
\section{Líneas de trabajo futuras}

Las dos principales líneas de trabajo serían estas:

\subsection{Guiado avanzado}
Como ya hemos expuesto en diferentes partes de esta memoria, la continuación mas inmediata de este proyecto sería, a partir de la trayectoria obtenida en dos dimensiones, y de conocer la posición del vehículo y la trayectoria, recrear virtualmente el entorno del vehículo para buscar una forma de guiado más avanzada. No solo ser capaz de guiarnos por lo que tenemos inmediatamente delante nuestro, sino ser capaces de anticiparnos a la trayectoria que tengamos seguir. 

Esto nos permitirá ajustar la velocidad, interpretar el radio de giro de una curva y estimar cuanto debería de girar el robot su sistema de dirección para cumplir ese radio de curvatura. Por lo que no solo sería más efectivo, sino que podría mejorar las velocidades de transporte.


\subsection{Reconocimiento de desniveles}
Otra mejora que se podría hacer es buscar una forma de recrear el entorno no solo en dos dimensiones, sino en tres, pudiendo reconocer objetos, obstáculos y desniveles. En este articulo\cite{ground_detection} podemos ver un pequeño proyecto que busca está funcionalidad, aunque está aplicado a exteriores, pero se podría aplicar a interiores.


\subsection{Robustez}
Más allá de estas dos funcionalidades, lo siguiente sería hacer el programa lo más robusto posible, de cara a poder filtrar todos los desperfectos que se generen en la imagen, siendo capaz de ver perfectamente la línea guía en todos los posibles entornos que se nos puedan plantear.


\bibliographystyle{plain}
\bibliography{bibliografia}

\end{document}
